% Generated draft: requires mathematical and source review.
\documentclass[11pt]{article}
\usepackage[margin=25mm]{geometry}
\usepackage{fontspec}
\setmainfont{Latin Modern Roman}
\usepackage{amsmath,amssymb,mathtools,unicode-math}
\setmathfont{Latin Modern Math}
\usepackage{enumitem,needspace,xurl,hyperref}
\setlength{\parindent}{0pt}
\setlength{\parskip}{5pt}
\setlength{\emergencystretch}{3em}
\newcommand{\N}{\mathbb{N}}
\newcommand{\Z}{\mathbb{Z}}
\newcommand{\Q}{\mathbb{Q}}
\newcommand{\R}{\mathbb{R}}
\newcommand{\C}{\mathbb{C}}
\newcommand{\F}{\mathbb{F}}
\newcommand{\A}{\mathbb{A}}
\newcommand{\PP}{\mathbb{P}}
\newcommand{\E}{\mathbb{E}}
\newcommand{\eps}{\varepsilon}
\newcommand{\dd}{\,\mathrm{d}}
\begin{document}
\section*{UnsolvedMath ID: 30002603}
\textbf{Source problem code:} catalogue entry
\subsection{Definitions and mathematical statement}
Let \(\{F(x):x\in\Z\}\) be an i.i.d. spatial field and let \(\{b(i):i\in\Z_{\geq0}\}\) be an i.i.d. temporal field, independent of the spatial field, with \(\E[b(i)]=0\). The two fields need not have the same marginal distribution. Define the product random potential by
\[
V(x,i)=F(x)b(i).
\]
The record additionally requires that ``the potential has a compactly supported density positive on the interior of its support.'' Read literally, this concerns the one-point marginal law of \(V(x,i)\); the supplied text does not resolve whether the intended density condition instead concerns \(F(x)\), \(b(i)\), or both. No alternative interpretation is imposed here.

For \(n\geq1\) and \(m\in\Z\) with \(|m|\leq n\), define the admissible point-to-point lazy paths, using a spatial origin as reference, by
\[
\mathcal P_n(0,m)=\left\{S:\{0,\ldots,n\}\to\Z:
S(0)=0,\ S(n)=m,\ |S(i+1)-S(i)|\leq1\quad(0\leq i<n)\right\}.
\]
If \(\mathcal A_n(S;V)\) denotes the action intended in the source report, write
\[
A_n(0,m)=\min_{S\in\mathcal P_n(0,m)}\mathcal A_n(S;V).
\]
Here \(\mathcal A_n\) is an unresolved part of the formulation: the supplied record does not give its formula, its time-summation convention, or any additional path cost. In particular, a specific sum of potential values is not asserted to be the source definition.

A conventional velocity-parametrized interpretation of the referenced shape function would be
\[
g(v)=\lim_{n\to\infty}\frac{A_n(0,\lfloor nv\rfloor)}{n},\qquad v\in[-1,1],
\]
provided the intended action and hypotheses yield this limit in the intended probabilistic sense. This display is a conditional interpretation, since the record does not specify the normalization, parametrization, mode of convergence, or whether the limiting function is deterministic.

\textbf{Question, subject to these unresolved definitions.} Under the stated assumptions on the product environment, does the minimum-action shape function have a nontrivial affine segment near velocity zero? A two-sided neighborhood formulation would ask whether there exist \(\delta\in(0,1]\) and \(\alpha,\beta\in\R\) such that
\[
g(v)=\alpha+\beta v\qquad\text{for every }v\in[-\delta,\delta].
\]
A one-sided formulation would instead require this identity on \([0,\delta]\) or \([-\delta,0]\). The supplied wording does not determine whether a two-sided affine segment is required, whether a one-sided segment suffices, or whether different affine pieces on the two sides of zero are allowed. Nor does it specify that the slope must vanish. These alternatives are not identified with one another.

The supplied dated curation labels the problem partially solved and reports a flat linear edge near the origin under an additional explicit power-law condition on a density at the edge of its support. That condition and the precise density to which it applies are not supplied. This reported result is therefore not a replacement for the question under the full extracted assumptions, and it is not independently verified here.
\subsection{Short English statement}
The problem asks whether the limiting minimum action for lazy lattice paths in an independent product environment is affine on a nontrivial interval near zero velocity. The supplied record leaves the action, the density assumption, the precise shape-function definition, and the required one-sided or two-sided interval unresolved. Supplied curation dated 2026-08-21 reports a partial result under an additional support-edge power-law hypothesis; this draft does not independently verify that assessment.
\subsection{Sources}
\begin{itemize}
\item Local UnsolvedMath registry, supplied complete source record: permanent ID \(30002603\), exact source problem code \texttt{catalogue entry}, title \emph{Linear Segments in Optimal-Path Shape Functions}. The record attributes the problem to the Oberwolfach Reports collection, gives proposed year 2014, and supplies no individual proposer.
\item Supplied report citation: \emph{Stochastic Analysis: Around the KPZ Universality Class} (2015), Oberwolfach Reports, pp.\ 1515--1568, DOI \texttt{10.4171/owr/2014/28}. Supplied link: \url{https://doi.org/10.4171/owr/2014/28}. The background describes Oberwolfach Reports as documenting mathematical workshops and their open-problem discussions.
\item Supplied literature source: Bakhtin--Khanin--Meszaros--Voltz, ``last passage percolation,'' \url{https://arxiv.org/abs/2310.08379}. This author attribution and abbreviated description are reproduced as supplied; no additional bibliographic details are inferred.
\item Supplied literature source: Oberwolfach Report 28/2014, \url{https://ems.press/content/serial-article-files/46518}.
\item Literature assessment checked at \texttt{2026-08-21}: supplied dated curation, not independently verified. It labels the problem partially solved and reports a flat linear edge near the origin for the product environment under an explicit power-law condition on the density at a support edge, covering a substantial regime but not every assumption in the extracted record. The supplied background explicitly calls for independent verification of time-sensitive status and cited claims.
\end{itemize}
\end{document}
